\chapter{Suicidal Ideation}

\section*{Definition}
Suicidal ideation encompasses thoughts of ending one's own life, ranging from passive wishes (\"I would be better off dead\") to active plans with intent, preparation, and means. It exists on a spectrum from occasional fleeting thoughts to persistent preoccupation with suicide, and must be distinguished from non-suicidal self-injury and death ideation unrelated to self-harm. Suicidal behavior includes preparatory acts, aborted attempts, and completed suicide.

\section*{What Happens in Your Body}
Suicidal behavior involves a complex interplay of neurobiological, psychological, and social factors. Dysfunction in the serotonergic system (reduced 5-HT1A binding, altered 5-HT2A receptors in prefrontal cortex) impairs behavioral inhibition and impulse control. The stress-diathesis model proposes that pre-existing vulnerabilities (genetic, epigenetic, early-life adversity) combined with acute stressors overwhelm coping mechanisms. Elevated cortisol and inflammatory markers, altered reward circuitry, and compromised decision-making in the ventromedial prefrontal cortex contribute to the transition from thoughts to action.

\section*{Causes}
\begin{itemize}
  \item Psychiatric disorders (most common): Major depressive disorder, Bipolar disorder (particularly depressive or mixed episodes), Schizophrenia (especially with command hallucinations), Borderline personality disorder, PTSD, Substance use disorders (especially alcohol, stimulants, opioids), Eating disorders (anorexia nervosa), Panic disorder, Body dysmorphic disorder
  \item Medical illness: Terminal or chronic painful conditions (cancer, HIV/AIDS, neurodegenerative disease), Traumatic brain injury, Epilepsy (especially temporal lobe), Chronic pain syndromes, Sleep disorders, TBI
  \item Psychosocial factors: Recent loss (relationship breakup, death, financial ruin), Social isolation or loneliness, Unemployment, Homelessness, Childhood trauma (abuse, neglect, ACEs), Legal problems, Perceived burdensomeness, Thwarted belongingness
  \item Genetic and biological: Family history of suicide, Altered serotonin metabolism, Epigenetic modifications at the SKA2 gene, Low cholesterol levels (association with violent suicide), HPA axis dysregulation
  \item Access to means: Firearms in the home, Access to lethal medications (opioids, tricyclic antidepressants, barbiturates), Access to pesticides in agricultural regions, High-risk environments (bridges, rooftops, subway)
\end{itemize}

\section*{When to See a Doctor}
See your doctor if:
\begin{itemize}
  \item Any expressed suicidal ideation, plan, or intent --- requires immediate evaluation regardless of setting
  \item Active intent, specific plan, or access to lethal means (highest acute risk)
  \item Recent suicide attempt or aborted attempt
  \item Suicidal thoughts accompanied by severe anxiety, agitation, psychosis, or substance intoxication
  \item New-onset or worsening suicidal ideation in you with known psychiatric illness
  \item Recent discharge from psychiatric hospitalization (highest risk period)
  \item Person who expresses hopelessness, feeling trapped, being a burden, or having no reason to live
  \item Acute change in mental status or new command hallucinations directing self-harm
\end{itemize}

\section*{Related Symptoms}
\begin{itemize}
  \item Depressed mood
  \item Anhedonia
  \item Hopelessness
  \item Worthlessness
  \item Anxiety and agitation
  \item Insomnia (especially early morning awakening)
  \item Social withdrawal
  \item Psychomotor retardation or agitation
  \item Impaired concentration
  \item Command hallucinations
  \item Substance use escalation
  \item Giving away possessions, saying goodbye (behavioral warning signs)
\end{itemize}
\textbf{Important:} Any expression of suicidal ideation warrants immediate clinical assessment; safety is the priority --- ask directly, ensure a safe environment, and never leave a suicidal patient alone.

\section*{Self-Care / Home Management}
\begin{itemize}
  \item Self-care is not appropriate for acute suicidal ideation --- immediate professional help is required (call or present to emergency services).
  \item For chronic passive ideation under professional care: maintain a safety plan developed with a therapist including warning signs, coping strategies, social supports, and emergency contacts.
  \item Remove or secure lethal means (firearms, medications, sharp objects) with the help of a trusted person.
  \item Avoid alcohol and non-prescribed drugs which disinhibit and worsen mood.
  \item Maintain sleep hygiene, basic nutrition, and social connection.
  \item National suicide prevention lifeline: 988 (US) or local crisis line stored in phone.
\end{itemize}

\section*{How Your Doctor Will Evaluate You}
\begin{itemize}
  \item Risk assessment is the cornerstone: inquire directly about suicidal thoughts, plan, intent, means, preparation, rehearsal, and past attempts (Columbia-Suicide Severity Rating Scale or Ask Suicide-Screening Questions).
  \item Assess protective factors (reasons for living, social support, treatment engagement) and risk factors (acute and chronic).
  \item Psychiatric evaluation for underlying diagnosis (MDD, bipolar, psychosis, substance use, personality disorder).
  \item Medical evaluation including toxicology screen and medical history to identify contributing conditions.
  \item Collateral history from family or friends whenever possible (patient may minimize risk).
  \item Safety planning with you using a collaborative, non-judgmental approach.
\end{itemize}

\section*{Treatment Options}
\begin{itemize}
  \item Imminent risk: Emergency psychiatric hospitalization for stabilization in a safe environment.
  \item Pharmacotherapy: Antidepressants (SSRIs/SNRIs with black box warning for increased suicidal thinking in initial weeks in people under 25), Mood stabilizers for bipolar disorder (lithium uniquely reduces suicide risk), Antipsychotics for psychotic disorders (clozapine reduces suicide risk in schizophrenia), Ketamine/esketamine for treatment-resistant depression with acute suicidal ideation.
  \item Psychotherapy: Cognitive behavioral therapy (CBT) for suicide prevention, Dialectical behavior therapy (DBT) for borderline personality disorder, Collaborative Assessment and Management of Suicidality (CAMS), Safety Planning Intervention.
  \item Means safety counseling: counsel patient and family to restrict access to firearms, medications, and other lethal methods.
  \item Close follow-up (within 24--72 hours of discharge) with contact frequency increased during high-risk periods.
  \item Treat comorbid substance use disorders as alcohol and drug use amplify suicide risk.
\end{itemize}

\section*{References}
\begin{thebibliography}{99}
\bibitem{suicidal_ideation:ref1} Bachmann S. "Epidemiology of suicide and the psychiatric perspective." \textit{International Journal of Environmental Research and Public Health}, 2018;15(7):1425.
\bibitem{suicidal_ideation:ref2} Kessler RC, Berglund P, Borges G, et al. "Trends in suicide ideation, plans, gestures, and attempts in the United States, 1990--1992 to 2001--2003." \textit{JAMA}, 2005;293(20):2487--2495.
\bibitem{suicidal_ideation:ref3} Mann JJ, Apter A, Bertolote J, et al. "Suicide prevention strategies: A systematic review." \textit{JAMA}, 2005;294(16):2064--2074.
\bibitem{suicidal_ideation:ref4} Zalsman G, Hawton K, Wasserman D, et al. "Suicide prevention strategies revisited: 10-year systematic review." \textit{The Lancet Psychiatry}, 2016;3(7):646--659.
\bibitem{suicidal_ideation:ref5} Turecki G, Brent DA, Gunnell D, et al. "Suicide and suicide risk." \textit{Nature Reviews Disease Primers}, 2019;5(1):73.
\bibitem{suicidal_ideation:ref6} Beck AT, Steer RA, Brown GK. \textit{Beck Depression Inventory-II Manual}. Psychological Corporation, 1996. [Suicidality assessment]
\end{thebibliography}
